\chapter{Envelopes and volumes of singularities}
\label{chap:envelopes}

The mass of a particular potential and the volume of its singularity type are different objects, even for smooth Hermitian potentials. We first construct the envelope operators, then introduce upper and lower volumes by varying the bounded part of a potential. The relative envelope class is defined without a mass bound. Its comparison with a class defined by equality of volumes will require the stated volume hypotheses.

\section{Obstacle envelopes}

A smooth real $(1,1)$-form $\theta$ is called \emph{big} if there are $\rho\in\PSH(X,\theta)$ and $\delta>0$ such that
\[
 \theta+\ddc\rho\geq\delta\omega.
\]
Here $\omega$ is any Hermitian metric. Changing that metric changes $\delta$, not the condition. Neither this definition nor the existence of $\theta$-plurisubharmonic functions requires $\theta$ to be closed.

For an obstacle $f$, write
\[
 P_\theta(f)=\left(\sup\{v\in\PSH(X,\theta):v\leq f
                      \ \hbox{q.e.}\}\right)^*,
 \qquad P_\theta(u,v)=P_\theta(\min(u,v)).
\]
The star denotes upper semicontinuous regularization, and the empty supremum is $-\infty$. Quasi-everywhere means outside a pluripolar set. This convention also gives meaning to differences of quasi-plurisubharmonic functions at common poles. For an upper semicontinuous obstacle, the pointwise definition gives the same competitors: a quasi-plurisubharmonic function lying below it almost everywhere lies below it everywhere, by the local submean inequality and upper semicontinuity of the obstacle.

We use the standard local envelope theorem as an external input: an upper envelope of plurisubharmonic functions, locally bounded above, has a plurisubharmonic upper semicontinuous regularization and differs from its unregularized supremum only on a pluripolar set. Applying local smooth corrections and the freezing argument of \cref{prop:bounded-products} gives the corresponding statement with the exact lower bound $\ddc v\geq-\theta$.

Here is a useful way to check the upper bound for an unbounded obstacle. If $f\leq A$ on a Borel set $E$ of positive $\mathrm{d}V$-measure and $h\leq f$ is a competitor, normalize $\widetilde h=h-\sup_Xh$. By \cref{thm:normalized-compactness},
\[
 (\sup_Xh-A)\operatorname{vol}_{\mathrm{d}V}(E)
 \leq\int_E|\widetilde h|\,\mathrm{d}V\leq C_\theta
 \quad\hbox{whenever }\sup_Xh>A.
\]
Thus the competitors have a common upper bound. In particular this applies to any measurable obstacle finite almost everywhere, by taking one of its finite sublevel sets of positive volume. This argument controls the supremum of the competitors; nonemptiness remains a separate requirement.

If $\PSH(X,\theta)$ is nonempty, its normalized envelope
\[
 V_\theta=P_\theta(0)
\]
is $\theta$-plurisubharmonic. Indeed, a normalized $\theta$-plurisubharmonic function is a competitor, and the family is bounded above by zero. Moreover,
\[
 u-\sup_Xu\leq V_\theta\qquad(u\in\PSH(X,\theta)).
\]
Thus $V_\theta$ has minimal singularities. If $\theta\geq0$, then $V_\theta=0$.

\begin{lemma}\label{lem:bounded-below-envelope}
Let $\eta\geq0$ be smooth. If an obstacle $f$ is finite almost everywhere and bounded below quasi-everywhere by a constant, then $P_\eta(f)$ is a bounded $\eta$-plurisubharmonic function.
\end{lemma}
\begin{proof}
The lower constant is a competitor. The preceding finite-sublevel argument gives a common upper bound for all competitors. The local envelope theorem therefore gives a bounded envelope.
\end{proof}

\begin{proposition}
\label{prop:continuous-obstacle-envelope}
If $f$ is continuous and $\omega$ is Hermitian, then $P_\omega(f)$ is continuous. Moreover,
\[
 \|P_\omega(f)-P_\omega(g)\|_\infty\leq\|f-g\|_\infty
 \quad(f,g\in C^0(X)).
\]
\end{proposition}
\begin{proof}
The constant $\min_Xf$ is a competitor, so $p=P_\omega(f)$ is bounded. The local envelope theorem gives $p\leq f$ quasi-everywhere, and continuity of $f$ makes this inequality pointwise. Choose smooth $v_j\downarrow p$ by \cref{thm:bounded-exact-regularization}. For every $\epsilon>0$, the continuous functions $v_j-f$ decrease to a function at most zero. Compactness shows that $v_j\leq f+\epsilon$ everywhere for all large $j$: otherwise their decreasing compact $\epsilon$-superlevel sets would have nonempty intersection. Hence $v_j-\epsilon$ is a competitor and
\[
 p\leq v_j\leq p+\epsilon.
\]
Thus $v_j\to p$ uniformly, proving continuity. The last inequality follows directly from the competitors and
$P_\omega(g+c)=P_\omega(g)+c$ for constants $c$.
\end{proof}

\begin{definition}\label{def:singularity-envelope}
For $u\in\PSH(X,\theta)$, set
\[
 P_\theta[u]
 =\left(\lim_{C\to+\infty}P_\theta(u+C,0)\right)^*
 =\left(\sup\{h\in\PSH(X,\theta):h\preceq u,\ h\leq0\}\right)^*.
\]
We call $u$ an \emph{envelope fixed point} if $P_\theta[u]=u$.
\end{definition}

The two formulas agree by taking the union of the competitor families. They also agree with the definition using $V_\theta$ in place of zero: every $\theta$-plurisubharmonic competitor below zero is below $V_\theta$. The envelope is a genuine potential, since the competitors contain $u-\sup_Xu$ and are bounded above by zero. In particular, $u\preceq P_\theta[u]$. The last relation is not an assertion that the two functions differ by a bounded amount.

Directly from the competitors, $P_\theta[u]$ is unchanged when $u$ is replaced by a potential of the same ordinary singularity type, and
\[
 u\preceq v\quad\Longrightarrow\quad P_\theta[u]\leq P_\theta[v].
\]
These statements concern ordinary singularities. Replacing $\preceq$ by the multiplier-ideal order would define a different operator, which will require a separate comparison.

\section{Contact sets}

A function is quasi-continuous if it becomes continuous after removal of open sets of arbitrarily small Bedford--Taylor capacity. Quasi-plurisubharmonic functions, and their differences away from common poles, have this property.

\Needspace{12\baselineskip}
\begin{proposition}\label{prop:smooth-obstacle-support}\index{Contact set!exponential approximation}
Let $\omega$ be Hermitian and $f\in C^\infty(X,\mathbb R)$. For $\beta>0$, denote by $u_\beta$ the smooth strictly $\omega$-plurisubharmonic solution of
\[
 \omega_{u_\beta}^n=\mathrm{e}^{\beta(u_\beta-f)}\omega^n.
\]
These solutions converge uniformly as $\beta\to+\infty$ to $p=P_\omega(f)$. Moreover,
\[
 \mathbf1_{\{p<f\}}\omega_p^n=0.
\]
\end{proposition}
\begin{proof}
Existence is \cref{thm:smooth-proper-existence}. Choose $A\geq0$ with $\ddc f\leq A\omega$. At a maximum of $u_\beta-f$,
\[
 \omega_{u_\beta}\leq\omega+\ddc f\leq(1+A)\omega.
\]
\Needspace{5\baselineskip}
The equation gives
\[
 u_\beta\leq f+\frac{C_f}{\beta},\qquad C_f=n\log(1+A).
\]
Thus $u_\beta-C_f/\beta$ is an envelope competitor, and
$u_\beta\leq p+C_f/\beta$.

The function $p$ is continuous by \cref{prop:continuous-obstacle-envelope}. Choose smooth strictly $\omega$-plurisubharmonic $v_j\downarrow p$ by \cref{thm:bounded-exact-regularization}; convergence is uniform by compactness. For any $\delta>0$, choose $j$ so that
\[
 p\leq v_j\leq p+\delta/2.
\]
Then $v_j-\delta\leq f$. Positivity and compactness give a number $c_j>0$ such that $\omega_{v_j}^n\geq c_j\omega^n$. Put $C_j=\max(0,-\log c_j)$. The function
\[
 w_{j,\beta}=v_j-\delta-C_j/\beta
\]
satisfies
\[
 \omega_{w_{j,\beta}}^n=\omega_{v_j}^n
 \geq \mathrm{e}^{-C_j}\omega^n
 \geq \mathrm{e}^{\beta(w_{j,\beta}-f)}\omega^n.
\]
Smooth exponential comparison gives $w_{j,\beta}\leq u_\beta$. Consequently
\[
 p-\delta-C_j/\beta\leq u_\beta\leq p+C_f/\beta.
\]
First let $\beta$ tend to infinity with $j$ fixed, so no uniform lower bound on $c_j$ is needed, and then let $\delta\downarrow0$. This proves uniform convergence. Bounded Bedford--Taylor continuity gives weak convergence of the measures.

On a compact subset of the open set $\{p<f\}$, uniform convergence gives $u_\beta-f\leq-\epsilon$ for all large $\beta$, with some $\epsilon>0$. There the measures have densities at most $\mathrm{e}^{-\beta\epsilon}$. Testing with continuous functions of compact support in that open set and passing to the bounded-product limit gives $\omega_p^n=0$ there.
\end{proof}

\begin{proposition}
\label{prop:lsc-obstacle-support}
If $\omega$ is Hermitian and $f$ is bounded and lower semicontinuous, then $P_\omega(f)$ is bounded, lies below $f$ quasi-everywhere, and its measure is concentrated on $\{P_\omega(f)=f\}$.
\end{proposition}
\begin{proof}
Suppose first that $f$ is continuous. Choose smooth $f_j\to f$ uniformly. The envelope Lipschitz bound gives $p_j=P_\omega(f_j)\to p=P_\omega(f)$ uniformly. On every compact subset of $\{p<f\}$, the inequalities $p_j<f_j$ hold on a neighborhood for all large $j$. The preceding proposition and bounded-product convergence then give $\omega_p^n=0$ on $\{p<f\}$.

For lower semicontinuous $f$, choose continuous $f_j\uparrow f$, with common upper and lower bounds. For example, for a fixed compatible distance $d$ one may take
\[
 f_j(x)=\inf_{y\in X}\bigl(f(y)+j d(x,y)\bigr).
\]
These functions are Lipschitz and increase to $f$: if convergence failed at $x$, approximate minimizers would approach $x$, contradicting lower semicontinuity.

The functions $p_j=P_\omega(f_j)$ increase, stay bounded, and have a bounded $\omega$-plurisubharmonic regularized limit $p_\infty=(\lim_jp_j)^*$. A countable union of pluripolar exceptional sets shows that $p_\infty\leq f$ quasi-everywhere. Thus $p_\infty\leq P_\omega(f)$.

For fixed $k$, the set $\{p_\infty<f_k\}$ is open. If $j\geq k$, then $p_j\leq p_\infty<f_k\leq f_j$ there, so the continuous-obstacle case gives zero measure on that open set. Increasing bounded Bedford--Taylor convergence applies after adding a fixed smooth local correction whose Hessian dominates $\omega$ and expanding against the resulting smooth background. Together with the weak-limit inequality for open sets, it yields
\[
 \omega_{p_\infty}^n(\{p_\infty<f_k\})=0.
\]
Taking the union over $k$ gives zero measure on $\{p_\infty<f\}$, hence also on $\{p_\infty<P_\omega(f)\}$ outside a pluripolar set. By bounded domination, \cref{thm:bounded-strict-domination}, $p_\infty\geq P_\omega(f)$. They are therefore equal, and the assertion follows.
\end{proof}

\begin{theorem}\label{thm:bounded-contact-support}
Let $\theta\geq0$ be smooth and let $f$ be bounded and quasi-continuous. Then $p=P_\theta(f)$ is bounded, $p\leq f$ quasi-everywhere, and
\[
 \mathbf1_{\{p<f\}}\mathrm{MA}_\theta(p)=0.
\]
\end{theorem}
\begin{proof}
First suppose $\theta=\omega$ is Hermitian. Choose $f_j\downarrow f$ quasi-everywhere as in \cref{lem:quasicont-lsc-approximation}. The bounded functions $p_j=P_\omega(f_j)$ decrease to $p$. Indeed their limit is a competitor for $p$, since it lies below $f$ quasi-everywhere; the reverse inequality follows from monotonicity.

For fixed $k$, put $H_k=\min\{1,(f-p_k)_+\}$. It is bounded and quasi-continuous. Where $H_k>0$ and $j\geq k$, we have $p_j\leq p_k<f\leq f_j$, outside a pluripolar set. The lower semicontinuous case therefore gives
\[
 \int_X H_k\,\omega_{p_j}^n=0.
\]
The weighted bounded-product convergence in \cref{sec:hermitian-capacity} gives
$\int_XH_k\,\omega_p^n=0$. The sets $\{p_k<f\}$ increase to $\{p<f\}$, proving support.

For general semipositive $\theta$, set $\theta_j=\theta+j^{-1}\omega$ with $\omega$ Hermitian. The functions $p_j=P_{\theta_j}(f)$ have the common bounds $\inf_Xf\leq p_j\leq\sup_Xf$ and decrease. Their limit is $\theta$-plurisubharmonic by distributional convergence, lies below $f$ quasi-everywhere, and hence equals $p=P_\theta(f)$. For each fixed $k$, the same coefficient $H_k=\min\{1,(f-p_k)_+\}$ gives
\[
 \int_XH_k\,\mathrm{MA}_{\theta_j}(p_j)=0\qquad(j\geq k).
\]
All $p_j$ are bounded potentials for the fixed Hermitian form $\theta+\omega$. Expand the products with $\theta_j$ against that fixed background. Bounded-product convergence and uniform capacity control pass the displayed identity to $\mathrm{MA}_\theta(p)$, including its quasi-continuous coefficient. Exhausting $\{p<f\}$ as before completes the proof.
\end{proof}


\Needspace{8\baselineskip}
\begin{proposition}
\label{prop:unbounded-contact-reduction}
Let $\theta\geq0$ be smooth. The bounded contact-support assertion of \cref{thm:bounded-contact-support} extends to every quasi-continuous obstacle $f$ for which $p=P_\theta(f)$ is a genuine potential. No upper mass bound is required.
\end{proposition}
\begin{proof}
The local envelope theorem and a countable Choquet family give $p\leq f$ outside a pluripolar set. Choose $M>\sup_Xp$ and replace $f$ by $\min(f,M)$. This does not change its envelope: $p$ itself is a competitor for the smaller obstacle. It suffices to prove support on the contact set for this smaller obstacle, because $p<M$ and therefore it has the same contact set as the original one.

Set $f_j=\max(f,-j)$ and $p_j=P_\theta(f_j)$. The constants $-j$ are admissible, so $p_j$ is bounded. \Cref{thm:bounded-contact-support} gives $p_j\leq f_j$ quasi-everywhere and
\[
 \mathbf1_{\{p_j<f_j\}}\mathrm{MA}_\theta(p_j)=0.
\]
The sequence $p_j$ decreases and stays above $p$. Its limit is a genuine $\theta$-plurisubharmonic function, lies below $f$ quasi-everywhere, and is therefore a competitor for $p$. Hence $p_j\downarrow p$.

Fix $k$ and $C>0$. Define the bounded nonnegative quasi-continuous function
\[
 H_{k,C}=\min\{1,(f-p_k)_+\}\,
                  \min\{1,(p+C)_+\},
 \qquad a_+=\max(a,0).
\]
Values on pluripolar exceptional sets are immaterial. Where this coefficient is positive, $p_k<f$ and $p>-C$. For $j\geq k$, we then have
$p_j\leq p_k<f\leq f_j$ and $p_j\geq p>-C$. Locality and \cref{thm:bounded-contact-support} imply
\[
 \int_X H_{k,C}\,
       \mathrm{MA}_\theta(\max(p_j,-C))=0.
\]
The truncated potentials are uniformly bounded and decrease to $\max(p,-C)$. Their measures converge weakly by bounded Bedford--Taylor continuity. They also have uniform local capacity bounds, so the same convergence holds against the fixed bounded quasi-continuous coefficient $H_{k,C}$: remove a set of small capacity, approximate on its complement by a continuous function, and then let that capacity tend to zero. Consequently,
\[
 \int_X H_{k,C}\,
       \mathrm{MA}_\theta(\max(p,-C))=0.
\]
By locality this says that the extended measure of $p$ vanishes on
$\{p_k<f\}\cap\{p>-C\}$. These sets increase to $\{p<f\}$ off the pole set when first $k$ and then $C$ tend to infinity. Continuity from below of positive measures proves the assertion, even if the total extended mass is infinite.
\end{proof}

\begin{theorem}
\label{thm:envelope-contact}
Let $\theta\geq0$ be smooth. If $f$ is quasi-continuous and $p=P_\theta(f)$ is a genuine potential, then $p\leq f$ quasi-everywhere and
\[
 \mathbf1_{\{p<f\}}\mathrm{MA}^{\mathrm{ext}}_\theta(p)=0.
\]
The total extended mass is allowed to be infinite.
\end{theorem}
\begin{proof}
The quasi-everywhere inequality follows from the local envelope theorem and Choquet extraction. The measure assertion is \cref{prop:unbounded-contact-reduction}.
\end{proof}

This proves the semipositive statement without a bigness or mass assumption. The corresponding conclusion in a semipositive big background is also \cite[Theorem 2.2]{ALS26}. Its proof uses precisely the bounded-to-unbounded passage handled above.

\begin{theorem}
\label{thm:envelope-contact-big}
Let $\theta$ be smooth and big. If $f$ is quasi-continuous, finite almost everywhere, and $p=P_\theta(f)$ is a genuine potential, then $p\leq f$ quasi-everywhere and
\[
\mathbf1_{\{p<f\}}\mathrm{MA}_\theta(p)=0.
\]
The extended measure may have infinite total mass. No bounded mass or positive lower-volume hypothesis is required.
\end{theorem}

The proof is given in \cref{thm:big-general-envelope-contact}, after establishing the local homogeneous replacement and continuous-obstacle case. A strict singular barrier repairs quasi-everywhere competitors, and signed bounded products control the final truncation argument. Under bounded mass, the contact conclusion is also \cite[Theorem 2.9 and Proposition 2.11]{BGL25}; that global mass hypothesis is not used in our proof. The ordered-contact inequality of \cref{prop:ordered-contact-general} applies to arbitrary smooth backgrounds without bigness or bounded mass.

\begin{proposition}\label{prop:rooftop-measure}
Under \cref{thm:envelope-contact} or \cref{thm:envelope-contact-big}, suppose $p=P_\theta(u,v)\in\PSH(X,\theta)$. Then
\[
 \mathrm{MA}_\theta(p)
 \leq\mathbf1_{\{p=u<v\}}\mathrm{MA}_\theta(u)
      +\mathbf1_{\{p=v\}}\mathrm{MA}_\theta(v).
\]
\end{proposition}
\begin{proof}
The envelope is below both obstacles; their upper semicontinuity gives this inequality everywhere. Its measure is supported on $\{p=\min(u,v)\}$ by the contact theorem. Partition this set into $\{p=u<v\}$ and $\{p=v\}$. On each part apply \cref{prop:ordered-contact-general}. All measures vanish on the pluripolar pole sets, which can be discarded throughout.
\end{proof}

\begin{proposition}\label{prop:fixed-point-finite}
Suppose $\theta$ is semipositive and big. Every envelope fixed point $\varphi$ satisfies
\[
 \mathrm{MA}^{\mathrm{ext}}_\theta(\varphi)
 \leq\mathbf1_{\{\varphi=0\}}\theta^n.
\]
In particular, its mass is finite, without assuming bounded mass on $X$.
\end{proposition}
\begin{proof}
We have $\varphi\leq0$. For every $j\geq0$, the envelope $P_\theta(\varphi+j,0)$ is both at least $\varphi$, since $\varphi$ is a competitor, and at most $P_\theta[\varphi]=\varphi$. Apply \cref{prop:rooftop-measure} to obtain
\[
 \mathrm{MA}^{\mathrm{ext}}_\theta(\varphi)
 \leq\mathbf1_{\{\varphi\leq-j\}}
          \mathrm{MA}^{\mathrm{ext}}_\theta(\varphi)
      +\mathbf1_{\{\varphi=0\}}\theta^n.
\]
One must not subtract potentially infinite measures here. Instead, restrict to $\{\varphi>-t\}$ and choose $j>t$. The first term vanishes. Letting $t\to+\infty$ proves the claim by the definition of the extended measure. This is the fixed-point argument in \cite[Proposition 2.12]{ALS26}.
\end{proof}

\section{Bounded domination in a semipositive background}

Call a smooth semipositive form $\theta$ \emph{non-collapsing on bounded potentials} if
\[
 \int_X\mathrm{MA}_\theta(w)>0
 \quad\hbox{for every }w\in\PSH(X,\theta)\cap L^\infty(X).
\]
The assertion is separate for each $w$. It is weaker than a uniform positive lower bound. Every Hermitian form has this property by \cref{cor:bounded-mass-positive}. Bigness also supplies it, by an envelope argument.

\begin{proposition}\label{prop:big-noncollapsing}
If a smooth real background $\theta$ is big, every bounded $u\in\PSH(X,\theta)$ has strictly positive Monge--Amp\`ere mass.
\end{proposition}
\begin{proof}
Choose $\rho\in\PSH(X,\theta)$ and $a>0$ with $\theta_\rho\geq a\omega$. Normalize $\rho\leq0$ and put
\[
 p=P_{a\omega}(u-\rho).
\]
The obstacle is finite almost everywhere and bounded below, so \cref{lem:bounded-below-envelope} gives bounded $p$. The obstacle is quasi-continuous off a pluripolar set. Contact support, applied in the Hermitian background $a\omega$, gives $\mathrm{MA}_{a\omega}(p)$ carried by
$D=\{p+\rho=u\}$, outside the pole set of $\rho$.

We have $p+\rho\leq u$ everywhere, by the canonical representatives, and
\[
 \theta+\ddc(p+\rho)
 =(a\omega+\ddc p)+(\theta_\rho-a\omega)
 \geq a\omega+\ddc p\geq0.
\]
On common finite truncation levels, positivity of the mixed products and ordered contact give
\[
 \mathbf1_D\mathrm{MA}_{a\omega}(p)
 \leq\mathbf1_D\mathrm{MA}^{\mathrm{ext}}_\theta(p+\rho)
 \leq\mathbf1_D\mathrm{MA}_\theta(u).
\]
Exhausting these levels proves the same inequalities for the extended measures. Since $p$ is bounded, its Hermitian Monge--Amp\`ere mass is positive by \cref{cor:bounded-mass-positive}. Integrating proves the assertion. No estimate uniform in $u$ is claimed.
\end{proof}

For semipositive big backgrounds, this gives the non-collapsing input needed below. The proof itself also applies to a big background which is not semipositive, whenever a bounded potential exists.


\begin{theorem}[Guedj--Lu]\label{thm:semipositive-bounded-domination}
Let $\theta\geq0$ be non-collapsing on bounded potentials. Suppose $u,v\in\PSH(X,\theta)\cap L^\infty(X)$ and $0\leq c<1$. If
\[
 \mathbf1_{\{u<v\}}\mathrm{MA}_\theta(u)
 \leq c\,\mathbf1_{\{u<v\}}\mathrm{MA}_\theta(v),
\]
then $u\geq v$.
\end{theorem}
\begin{proof}
We give the envelope argument of \cite[Proposition 2.8]{GL22}. If the conclusion failed, choose $a>0$ such that $E=\{u<v-a\}$ is nonempty. This set has positive volume by \cref{lem:strict-set-positive-volume}. Fix $C$ with $|u|,|v|\leq C$.

For $b>1$ define the bounded envelope
\[
 w_b=P_\theta\bigl(bu-(b-1)v\bigr),\qquad
 D_b=\{w_b=bu-(b-1)v\}.
\]
Constants bounded below are competitors, since $\theta\geq0$. Bounded contact support gives $\mathrm{MA}_\theta(w_b)$ concentrated on $D_b$. The potential
$b^{-1}w_b+(1-b^{-1})v$ lies below $u$ quasi-everywhere, hence everywhere by the canonical representatives, and equals $u$ on $D_b$ outside a pluripolar set. Ordered contact and positivity of the mixed terms give
\[
 \mathbf1_{D_b}\left[b^{-n}\mathrm{MA}_\theta(w_b)
              +(1-b^{-1})^n\mathrm{MA}_\theta(v)\right]
 \leq\mathbf1_{D_b}\mathrm{MA}_\theta(u).
\]
Choose $b$ sufficiently large that $(1-b^{-1})^n>c$. On $D_b\cap\{u<v\}$ this forces $\mathrm{MA}_\theta(w_b)=0$. Non-collapsing gives positive total mass for $w_b$, so some mass is carried by $D_b\cap\{u\geq v\}$. At a point of that set outside the exceptional pluripolar set,
\[
 w_b=bu-(b-1)v=u+(b-1)(u-v)\geq-C.
\]
Consequently $\sup_Xw_b\geq-C$, uniformly as $b\to\infty$.

On the fixed set $E$, the obstacle bound gives
$w_b\leq C-(b-1)a$ quasi-everywhere. Hence
\[
 \int_X|w_b-\sup_Xw_b|\,\mathrm{d}V
 \geq\bigl((b-1)a-2C\bigr)\operatorname{vol}_{\mathrm{d}V}(E)
\]
for all sufficiently large $b$. This contradicts the normalized $L^1$ bound of \cref{thm:normalized-compactness}. No uniform estimate for the masses of $w_b$ was used.
\end{proof}

The same argument as in \cref{cor:bounded-exponential-comparison} now gives exponential comparison for bounded potentials in every such semipositive background. Removing boundedness requires further control: an envelope used in the proof may then have infinite mass or fail to lie in the desired singularity class.

\section{Upper and lower volumes}

In this section $X$ has the bounded mass property and $\theta$ is big. All non-pluripolar measures below are finite. The extended definitions also make sense for semipositive backgrounds without bounded mass, but then the upper volume can be infinite.

\begin{definition}\label{def:singularity-volumes}
For $u\in\PSH(X,\theta)$, define
\[
 \begin{split}
 \underline{\operatorname{vol}}_\theta(u)
  &=\inf_{v\in\PSH(X,\theta),\ v\simeq u}
       \int_X\mathrm{MA}_\theta(v),\\
 \overline{\operatorname{vol}}_\theta(u)
  &=\sup_{v\in\PSH(X,\theta),\ v\simeq u}
       \int_X\mathrm{MA}_\theta(v).
 \end{split}
\]
These are the lower and upper volumes of the singularity type of $u$ in the fixed background. We write $\underline{\operatorname{vol}}(\theta)$ and $\overline{\operatorname{vol}}(\theta)$ for their values at $V_\theta$.
\end{definition}

This is \cite[Definition 3.1]{BGL25}, expressed in potentials. Clearly,
\[
 0\leq\underline{\operatorname{vol}}_\theta(u)
 \leq\int_X\mathrm{MA}_\theta(u)
 \leq\overline{\operatorname{vol}}_\theta(u)<+\infty.
\]
The uniform finiteness at the right is the content of the bounded mass estimates, not merely the finiteness of each individual integral.

The two volumes are invariant under bounded changes of $u$ by definition. They are also invariant under the smooth gauge change $(\theta,u)\mapsto(\theta+\ddc f,u-f)$: the map $v\mapsto v-f$ identifies the competitor currents and their measures. Scaling gives
\[
 \underline{\operatorname{vol}}_{t\theta}(tu)
  =t^n\underline{\operatorname{vol}}_\theta(u),\qquad
 \overline{\operatorname{vol}}_{t\theta}(tu)
  =t^n\overline{\operatorname{vol}}_\theta(u)\quad(t>0).
\]
If $\theta$ is closed, both volumes equal $\int_X\mathrm{MA}_\theta(u)$; this uses invariance of the non-pluripolar mass under equivalent singularities, recalled in \cite[Example 3.2]{BGL25}. For a non-closed background this equality can fail even at $u=0$, as \cref{ex:torus-mass} shows.

\Needspace{12\baselineskip}
\begin{theorem}\label{thm:volume-monotonicity}
Suppose $\theta_1$ is big, $\theta_1\leq\theta_2$, and $u_i\in\PSH(X,\theta_i)$ satisfy $u_1\preceq u_2$. On a manifold with bounded mass,
\[
 \underline{\operatorname{vol}}_{\theta_1}(u_1)
 \leq\underline{\operatorname{vol}}_{\theta_2}(u_2),\qquad
 \overline{\operatorname{vol}}_{\theta_1}(u_1)
 \leq\overline{\operatorname{vol}}_{\theta_2}(u_2).
\]
\end{theorem}
\begin{proof}
We reconstruct \cite[Proposition 3.7]{BGL25} in the big regime. Shift $u_1$ so that $u_1\leq u_2$. Choose $v_2\simeq u_2$ in $\PSH(X,\theta_2)$ and set
\[
 p_t=P_{\theta_1}(u_1+t,v_2).
\]
For $t>0$, a sufficiently low constant translate of $u_1$ is a competitor; also $p_t\leq u_1+t$. Thus $p_t\simeq u_1$. The contact theorem, applied in each background on the corresponding contact part, and positivity of the mixed products give
\[
 \mathrm{MA}_{\theta_1}(p_t)
 \leq\mathbf1_{\{u_1\leq v_2-t\}}\mathrm{MA}_{\theta_1}(u_1)
      +\mathrm{MA}_{\theta_2}(v_2).
\]
Indeed $p_t$ is also $\theta_2$-plurisubharmonic and is below $v_2$, so the contact inequality in that background applies. The first measure on the right has mass tending to zero: its indicators tend to zero off the pole sets, and the measure is finite and non-pluripolar. Integrate, let $t\to+\infty$, and then take the infimum over $v_2$.

For the upper volume, choose any $v_1\simeq u_1$ in $\PSH(X,\theta_1)$ and put $w_t=\max(v_1,u_2-t)$. It is $\theta_2$-plurisubharmonic and $w_t\simeq u_2$. On $D_t=\{v_1>u_2-t\}$, locality and mixed-product positivity imply
\[
 \mathbf1_{D_t}\mathrm{MA}_{\theta_1}(v_1)
 \leq\mathbf1_{D_t}\mathrm{MA}_{\theta_2}(w_t).
\]
The sets $D_t$ exhaust the complement of a pluripolar set. Hence integration and passage to the limit give
$\int_X\mathrm{MA}_{\theta_1}(v_1)\leq
\overline{\operatorname{vol}}_{\theta_2}(u_2)$.
Take the supremum over $v_1$. This proof does not require $v_1$ to be bounded.
\end{proof}

The theorem orders the backgrounds and the potentials separately. These hypotheses should not be confused with ordering the two currents $\theta_i+\ddc u_i$.

\begin{proposition}\label{prop:volume-addition}
Assume bounded mass. Let $\theta$ be big, $u\in\PSH(X,\theta)$, and $v\in\PSH(X,\eta)$ for a smooth real form $\eta$. Then, for either the lower or the upper volume,
\[
 \operatorname{vol}_\theta(u)
 \leq\operatorname{vol}_{\theta+\eta}(u+v).
\]
\end{proposition}
\begin{proof}
The form $\theta+\eta$ is big, by adding the positive current $\eta+\ddc v$ to a strictly positive current in the class of $\theta$. For the upper volume, add $v$ to each competitor for $u$. Positivity of mixed products shows that the resulting mass is at least the original mass.

For the lower volume, choose $w\simeq u+v$ in $\PSH(X,\theta+\eta)$, and let $p=P_\theta(w-v)$. Since $w-v=u+O(1)$ outside the pole sets, the envelope is genuine and $p\simeq u$. The contact theorem says that its measure is carried by $D=\{p+v=w\}$. Furthermore $p+v\leq w$ everywhere, since this inequality holds off a pluripolar set between canonical quasi-plurisubharmonic functions. Thus
\[
 \mathrm{MA}_\theta(p)
 \leq\mathbf1_D\mathrm{MA}_{\theta+\eta}(p+v)
 \leq\mathbf1_D\mathrm{MA}_{\theta+\eta}(w).
\]
Integrating and taking the infimum over $w$ proves the assertion. This is the envelope proof of \cite[Proposition 2.5]{XL26}.
\end{proof}

\section{Relative full mass through convex combinations}

For $u,\varphi\in\PSH(X,\theta)$ and $b\geq1$, define the competitor family
\[
 \mathcal A_b(u,\varphi)=
 \left\{h\in\PSH(X,\theta):
 b^{-1}h+(1-b^{-1})\varphi\leq u\right\}.
\]
For $b=1$ the inequality means $h\leq u$. For $b>1$ the displayed sum of quasi-plurisubharmonic functions is defined also at poles. This notation avoids assigning a value to $bu-(b-1)\varphi$ at a common pole. When the family is nonempty, its upper envelope is denoted
\[
 Q_b(u,\varphi)=P_\theta(bu-(b-1)\varphi).
\]
For big $\theta$, the family is uniformly bounded above. Indeed, choose $m$ such that $D=\{\varphi>-m\}$ has positive smooth volume $d$. On $D$, every competitor satisfies $h\leq L:=b\sup_Xu+(b-1)m$. Choose $C>0$ with $\theta\leq C\omega$. The normalized $L^1$ estimate \eqref{eq:normalized-lone}, applied after scaling, gives $\int_X|h-\sup_Xh|\,\mathrm{d}V\leq K$. Thus $\sup_Xh\leq L+K/d$ whenever $\sup_Xh>L$, and the bound is immediate otherwise. The local envelope theorem then shows that $Q_b$ is a genuine potential. The defining convex inequality holds everywhere by the canonical quasi-plurisubharmonic representatives. These facts are also recorded in \cite[Propositions 2.4 and Lemma 2.5]{ALS26} and \cite[Section 2.1]{XL26}.

\begin{definition}\label{def:relative-envelope-class}
Define
\[
 \mathcal E_{\mathrm{env}}(X,\theta,\varphi)
 =\{u\in\PSH(X,\theta):u\preceq\varphi,
       \ \mathcal A_b(u,\varphi)\ne\varnothing
       \text{ for every }b\geq1\}.
\]
Under bounded mass we also define
\[
 \mathcal E_{\mathrm{vol}}(X,\theta,\varphi)
 =\{u\in\PSH(X,\theta):u\preceq\varphi,
 \ \overline{\operatorname{vol}}_\theta(u)
   =\overline{\operatorname{vol}}_\theta(\varphi)\}.
\]
\end{definition}

The first is the definition used in \cite[Definition 3.2]{ALS26} for semipositive big $\theta$; the second is \cite[Definition 3.2]{XL26}. Our temporary subscripts keep their definitions visible before proving their equivalence in a common regime.

\begin{proposition}\label{prop:relative-elementary}
The class $\mathcal E_{\mathrm{env}}(X,\theta,\varphi)$ is convex, contains all potentials equivalent to $\varphi$, and is stable under bounded changes of its members. If $u$ belongs to it, $u\leq v$ and $v\preceq\varphi$, then $v$ belongs to it. In particular, it is stable under taking maxima of two members.
\end{proposition}
\begin{proof}
If $u\geq\varphi-C$, use $h=\varphi-bC$ for each $b$. More generally, if $u'\geq u-C$ and $h\in\mathcal A_b(u,\varphi)$, then $h-bC\in\mathcal A_b(u',\varphi)$. The singularity-order condition persists under bounded changes. Convex combinations of two competitors give competitors for the same convex combination of their target functions. Finally a competitor for $u$ is a competitor for any larger $v$; the assumed relation $v\preceq\varphi$ supplies the other defining condition.
\end{proof}

\begin{theorem}\label{thm:relative-equivalence}
Suppose $X$ has bounded mass, $\theta$ is big, and
$\underline{\operatorname{vol}}_\theta(\varphi)>0$. Then
\[
 \mathcal E_{\mathrm{env}}(X,\theta,\varphi)
 =\mathcal E_{\mathrm{vol}}(X,\theta,\varphi).
\]
Every member also satisfies
$\underline{\operatorname{vol}}_\theta(u)
=\underline{\operatorname{vol}}_\theta(\varphi)$.
\end{theorem}
\begin{proof}
First assume the envelope condition. For each $a\in(0,1)$ there is $h\in\PSH(X,\theta)$ such that $a\varphi+(1-a)h\leq u$. By \cref{thm:volume-monotonicity,prop:volume-addition} and homogeneity, for either volume,
\[
 a^n\operatorname{vol}_\theta(\varphi)
 \leq\operatorname{vol}_\theta(a\varphi+(1-a)h)
 \leq\operatorname{vol}_\theta(u)
 \leq\operatorname{vol}_\theta(\varphi).
\]
Let $a\to1$. This proves equality of both volumes in this direction, even without positive lower volume.

Conversely, suppose equality of the upper volumes, and fix $b>1$. Shift the potentials so that $u\leq\varphi\leq0$, and put
\[
 u_j=\max(u,\varphi-j),\qquad q_j=Q_b(u_j,\varphi).
\]
Each $q_j$ is genuine and equivalent to $\varphi$: the lower competitor $\varphi-bj$ and the defining inequality give
$\varphi-bj\leq q_j\leq\varphi$.
The sequence decreases. If its limit were $-\infty$, the compactness and Hartogs lemma for normalized quasi-plurisubharmonic functions would give $\sup_Xq_j\to-\infty$.

Set $C_j=\{b^{-1}q_j+(1-b^{-1})\varphi=u_j\}$. The contact theorem and positivity of mixed products give
\[
 \mathrm{MA}_\theta(q_j)
 \leq b^n\mathbf1_{C_j}\mathrm{MA}_\theta(u_j).
\]
On $C_j\cap\{q_j\leq-bk\}$, the inequality $\varphi\leq0$ implies $u_j\leq-k$, hence $u\leq-k$. For $j>k$, locality on $\{u>-k\}\subset\{u>\varphi-j\}$ therefore yields
\[
 \begin{split}
 \int_{\{q_j\leq-bk\}}\mathrm{MA}_\theta(q_j)
 &\leq b^n\int_{\{u\leq-k\}}\mathrm{MA}_\theta(u_j)\\
 &\leq b^n\left(
 \overline{\operatorname{vol}}_\theta(\varphi)
 -\int_{\{u>-k\}}\mathrm{MA}_\theta(u)\right).
 \end{split}
\]
If $\sup_Xq_j\to-\infty$, the domain on the left is all of $X$ for large $j$. It follows, on letting $k\to+\infty$, that
\[
 \underline{\operatorname{vol}}_\theta(\varphi)
 \leq b^n\left(
 \overline{\operatorname{vol}}_\theta(\varphi)
 -\int_X\mathrm{MA}_\theta(u)\right).
\]
The same collapse would hold with any $u'\simeq u$ in place of $u$: a competitor for either target gives one for the other by a constant translation. Repeating the argument and taking the supremum of $\int_X\mathrm{MA}_\theta(u')$ gives
\[
 \underline{\operatorname{vol}}_\theta(\varphi)
 \leq b^n\left(
 \overline{\operatorname{vol}}_\theta(\varphi)
 -\overline{\operatorname{vol}}_\theta(u)\right)=0,
\]
a contradiction. Thus $Q_b(u,\varphi)$ is genuine for every $b$. The argument is the full-volume case of \cite[Lemma 3.1]{XL26}; the proof uses bounded mass and the displayed positive lower volume, without a separate global positive volume hypothesis.
\end{proof}

\section{Which model condition?}

\begin{proposition}\label{prop:relative-envelope-equality}
For a big background, if $u\in\mathcal E_{\mathrm{env}}(X,\theta,\varphi)$, then
\[
 P_\theta[u]=P_\theta[\varphi].
\]
This implication does not require bounded mass.
\end{proposition}
\begin{proof}
The relation $u\preceq\varphi$ gives one inequality. For $a\in(0,1)$ choose $h_a$ with $a\varphi+(1-a)h_a\leq u$. Combining normalized competitors for $P_\theta[\varphi]$ and $P_\theta[h_a]$ gives
\[
 aP_\theta[\varphi]+(1-a)P_\theta[h_a]\leq P_\theta[u].
\]
Indeed, their convex combination is below zero and is more singular than $u$; taking the two upper envelopes preserves the inequality by local regularization. For every potential $h$, $P_\theta[h]$ has supremum zero, because it lies above $h-\sup_Xh$ and below zero. The normalized $L^1$ bound therefore applies uniformly to $P_\theta[h_a]$. Along $a\to1$ the left-hand side tends to $P_\theta[\varphi]$ in $L^1$. The inequality passes to the limit almost everywhere and then everywhere for the canonical representatives. This proves the other inequality.
\end{proof}

Two additional positivity conditions accompany envelope fixed points in the literature. We state them separately.

For semipositive big $\theta$, the condition used by Alehyane--Lu--Salouf is
\begin{equation}\label{eq:als-model-condition}
 P_\theta[\varphi]=\varphi,\qquad
 P_{\varepsilon\theta}(\varphi)\not\equiv-\infty
 \quad\text{for some }0<\varepsilon<1.
\end{equation}
The second requirement leaves room to realize the singularity below a smaller background. For a big, possibly non-semipositive form under bounded mass, Li uses
\begin{equation}\label{eq:li-model-condition}
 P_\theta[\varphi]=\varphi,\qquad
 \underline{\operatorname{vol}}_\theta(\varphi)>0.
\end{equation}
We refer to the formulas when invoking their theorems, rather than silently treating the two meanings of ``model'' as interchangeable.

In the semipositive regime, \cref{prop:fixed-point-finite} supplies finiteness of the mass in \eqref{eq:als-model-condition}. The smaller-background condition further implies positivity of the mass of every potential equivalent to $\varphi$, by \cite[Proposition 2.10 and the discussion after Definition 2.13]{ALS26}. This is individual positivity; a positive infimum over all such potentials is a stronger assertion. Under bounded mass, positive lower volume implies the smaller-background condition by \cite[Proposition 2.11]{ALS26}. That paper does not assume bounded mass for its same-singularity-type equations with $L^p$ data.

The model conditions also clarify a limitation of raw mass. On a manifold without bounded mass, \cref{thm:infinite-mass-witness} gives potentials with infinite mass and trivial multiplier ideals at every weight. Meanwhile every envelope fixed point for a semipositive big background has finite mass. Thus multiplier ideals alone do not control the unnormalized mass of a representative. Any extension of the source book's $\mathcal I$-goodness equivalences must account for the envelope and the chosen volume invariant.

\section{Infinite mass inside the envelope class}

The same series construction gives a stronger obstruction. It concerns the relative class itself, with the smooth reference potential zero.

\begin{proposition}\label{prop:infinite-envelope-class}
If $v_+(\omega)=+\infty$, the potential $U$ in \cref{thm:infinite-mass-witness} can be chosen so that, in addition,
\[
 U\in\mathcal E_{\mathrm{env}}(X,\omega,0),\qquad P_\omega[U]=0.
\]
The same conclusions hold for $\varepsilon U$, $0<\varepsilon\leq1$. Each of these potentials has infinite extended mass and trivial multiplier ideals at every positive weight.
\end{proposition}
\begin{proof}
Use exactly the series in that theorem, and write
\[
 U=A_N+V_N,\qquad
 A_N=\sum_{j\leq N}2^{-j}u_j,\qquad
 V_N=\sum_{j>N}2^{-j}u_j,\qquad \delta_N=2^{-N}.
\]
The head $A_N$ is smooth. The tail is nonpositive and belongs to $\PSH(X,\delta_N\omega)$, by taking decreasing limits of its finite sums. It is therefore also $\omega$-plurisubharmonic. Since $V_N-U=-A_N$ is smooth on compact $X$, we have $V_N\simeq U$ for every $N$. Each $V_N$ is a competitor for $P_\omega[U]$, so
\[
 V_N\leq P_\omega[U]\leq0.
\]
The tails increase to zero wherever the series for $U$ is finite, hence almost everywhere. The envelope consequently equals zero almost everywhere, and thus everywhere by its canonical quasi-plurisubharmonic representative.

Fix $b>1$ and choose $N$ so large that $b\delta_N\leq1$. Then
\[
 h_b=bV_N+\min_X(bA_N)
\]
is $\omega$-plurisubharmonic: its curvature is bounded below by $-b\delta_N\omega\geq-\omega$. Moreover $h_b\leq bU$. This gives a member of $\mathcal A_b(U,0)$ for every $b$, proving the relative-class assertion. For $\varepsilon U$, use the tails $\varepsilon V_N$ and choose $N$ with $b\varepsilon\delta_N\leq1$. The infinite-mass and integrability conclusions are already proved in \cref{cor:bmp-characterization}.
\end{proof}

\begin{corollary}\label{cor:bmp-relative-finiteness}
The following are equivalent for a Hermitian metric $\omega$:
\begin{enumerate}
\item $v_+(\omega)<+\infty$;
\item every member of $\mathcal E_{\mathrm{env}}(X,\omega,0)$ has finite extended mass;
\item every $u\in\PSH(X,\omega)$ with $P_\omega[u]=0$ has finite extended mass.
\end{enumerate}
\end{corollary}
\begin{proof}
Bounded mass gives finiteness for all $\omega$-plurisubharmonic functions by \cref{prop:bmp-finite}. If it fails, \cref{prop:infinite-envelope-class} contradicts each of the other two assertions.
\end{proof}

Thus the name ``relative full mass'' for the envelope-defined class does not guarantee finite total mass on an arbitrary Hermitian manifold. There is no conflict with the prescribed-equation theorems: their solutions have finite right-hand sides, and their convergence arguments impose additional control. Nor is there a conflict with \cref{prop:fixed-point-finite}: $P_\omega[U]=0$ says that the envelope of $U$ is smooth, not that $U$ itself is an envelope fixed point. This distinction is essential in interpreting modelization outside bounded mass.
